% ============================================================
% 05-results.tex — V13 Results (condensed)
% ============================================================
\section{Results}
\label{sec:results}

%------------------------------------------------
\subsection{Scenario 1 --- Architecture-Level Comparison}
\label{sec:s1_results}

Table~\ref{tab:arch_comparison} compares the three configurations on the 100-query held-out benchmark. Because pre-generation evidence retrieval is held fixed across configurations, Source Recall (78.6\%) and Source AP (70.8\%) are identical across all three architectures. The observed differences in final answer quality and operational cost may be associated with the combined orchestration changes described in Section~\ref{sec:exp_setup}.

\begin{table}[!htbp]
\centering
\caption{Scenario~1: Architecture-level comparison ($N=100$). $\text{Pass}_{.50}$ = Automated Fact Pass Rate ($\ge 50\%$ required facts). Tools/Q = mean tool invocations ($\downarrow$). In/Out Tok = specialist agent tokens. Bold = best reliability / lowest cost.}
\label{tab:arch_comparison}
{\small\setlength{\tabcolsep}{3.5pt}
\begin{tabular}{lcccccccc}
\toprule
\textbf{Arch.} & \textbf{Pass$_{.50}$} & \textbf{Fact} & \textbf{Src.R} & \textbf{Src.AP} & \textbf{Tools/Q} & \textbf{In Tok} & \textbf{Out Tok} & \textbf{Lat.\,(ms)} \\
\midrule
Single Agent   & 66.0 & 57.5 & 78.6 & 70.8 & 0.41 & 6296 & 1530 & \textbf{5159} \\
Routed Gen.    & 66.0 & 54.2 & 78.6 & 70.8 & 0.31 & \textbf{3473} & 1415 & 8921 \\
\textbf{CTU-Chat} & \textbf{76.0} & \textbf{61.0} & 78.6 & 70.8 & \textbf{0.28} & 3690 & \textbf{2092} & 10652 \\
\bottomrule
\end{tabular}}
\end{table}

\textbf{Pairwise bootstrap analysis} (10,000 resamples):
(1)~CTU-Chat vs.\ Single Agent: CTU-Chat achieves a 10.0\,pp higher $\text{Pass}_{\ge 50\%}$ point estimate (76.0\% vs.\ 66.0\%, paired 95\%\,bootstrap CI $[+2.0, +19.0]$\,pp; exact McNemar $p{=}0.0414$), alongside higher Fact Coverage ($\Delta{=}+3.58$\,pp, CI $[-1.2, +8.5]$\,pp) and 41.4\% fewer specialist input tokens, under differing measurement boundaries (10{,}652\,ms end-to-end pipeline vs.\ 5{,}159\,ms post-retrieval generation).
(2)~CTU-Chat vs.\ Routed Generic: specialist policies raise Fact Coverage by $+6.85$\,pp (CI $[+0.7, +13.0]$\,pp, $p{<}0.05$).
(3)~Routed Generic vs.\ Single Agent: both yield $66.0\%$ pass rate, but Routed Generic records lower Fact Coverage ($54.2\%$ vs.\ $57.5\%$). The Single Agent suffers prompt interference from all 11 tool schemas ($0.41$ tools/query), while Routed Generic restricts visibility but lacks domain-specific reasoning directives---demonstrating that tool scoping and specialist prompting must operate conjointly.

On the 20 cross-domain queries, \system{} achieves an 80.0\% pass rate (vs.\ 70.0\% Single, 65.0\% Generic). Evaluating sensitivity across thresholds, \system{} demonstrates a higher automated pass rate at the baseline 50\% fact-coverage threshold (unadjusted exact McNemar $p{=}0.0414$). This margin narrows at 60\% (54.0\% vs.\ 50.0\%) and reverses at stricter thresholds (e.g., 27.0\% vs.\ 30.0\% at 75\%). However, the difference at 75\% remains statistically uncertain ($p{=}0.5488$, 95\% CI $[-10.0, +3.0]$\,pp), indicating neither architecture holds a definitive advantage under stringent fact-matching criteria.

\textbf{Operational Trade-offs:}
Table~\ref{tab:arch_comparison} reveals accuracy--efficiency trade-offs. \system{} achieves the highest reliability ($\text{Pass}_{.50} = 76.0\%$) with $0.28$ tools/query, but its end-to-end pipeline latency ($10{,}652$\,ms, including routing and contract verification) and Single Agent post-retrieval timing ($5{,}159$\,ms) span different boundaries, precluding direct speed comparison. Within multi-agent variants, specialist prompts add $+1{,}731$\,ms and $+677$ output tokens over generic workers for domain-tailored reasoning. Specialist scoping reduces recorded input tokens ($3{,}690$ vs.\ $6{,}296$), offset by ${\sim}650$ upstream routing tokens per query. Single-agent designs suit single-domain FAQ workflows with small tool pools; \system{} is preferable where heterogeneous evidence and deterministic arithmetic carry administrative consequences.

%------------------------------------------------
\subsection{Scenario 2 --- Architectural Ablation}
\label{sec:s2_results}

Table~\ref{tab:ablation} removes one mechanism at a time from Full CTU-Chat ($N{=}100$). Deltas are relative to the full system.

\begin{table}[!htbp]
\centering
\caption{Scenario~2: Architectural ablation results ($N=100$). Panel~A reports overall benchmark differences ($\Delta$ relative to Full \system{}). Panel~B reports Automated Fact Pass Rate ($\text{Pass}_{\ge 50\%}$, \%) across query families.}
\label{tab:ablation}
{\small\setlength{\tabcolsep}{3.5pt}
\begin{tabular}{lccccc}
\toprule
\multicolumn{6}{l}{\textbf{Panel A: Overall Benchmark Ablations}} \\
\textbf{Ablated Component} & \textbf{$\Delta$Pass$_{.50}$} & \textbf{$\Delta$Fact} & \textbf{$\Delta$Src.R} & \textbf{$\Delta$Tok} & \textbf{$\Delta$Lat} \\
\midrule
No Spec.~Policy  & $-$10.0 & $-$6.85 & 0.0  & $-$217  & $-$1731 \\
Full Tool Vis.   & 0.0     & +0.93   & 0.0  & +1270   & +296 \\
Uniform Text-RAG & $-$9.0  & $-$2.18 & $-$12.3 & $-$300  & $-$1910 \\
No Route Rep.    & $-$2.0  & $-$1.27 & 0.0  & $-$219  & $-$1610 \\
No Determ.~Calc.$^\dagger$ & $-$4.0 & $-$0.93 & 0.0 & $-$782 & $-$4540 \\
\midrule
\multicolumn{6}{l}{\textbf{Panel B: Fact-Grounded Pass Rate ($\text{Pass}_{\ge 50\%}$) by Query Family}} \\
\textbf{Configuration} & \textbf{Graph} ($n{=}19$) & \textbf{Struct-fin} ($n{=}28$) & \textbf{Narr-reg} ($n{=}33$) & \multicolumn{2}{c}{\textbf{Cross} ($n{=}20$)} \\
\midrule
\textbf{Full CTU-Chat} & 63.2 & \textbf{75.0} & 81.8 & \multicolumn{2}{c}{\textbf{80.0}} \\
No Spec.~Policy  & 52.6 & 67.9 & 72.7 & \multicolumn{2}{c}{65.0} \\
Full Tool Vis.   & \textbf{68.4} & 67.9 & \textbf{84.8} & \multicolumn{2}{c}{\textbf{80.0}} \\
Uniform Text-RAG & 63.2 & 60.7 & 78.8 & \multicolumn{2}{c}{60.0} \\
No Route Rep.    & \textbf{68.4} & \textbf{78.6} & 78.8 & \multicolumn{2}{c}{65.0} \\
\bottomrule
\multicolumn{6}{l}{\rule{0pt}{2.5ex}\footnotesize $^\dagger$Full benchmark with calculators disabled. $\Delta$Tok/Lat in tokens/ms per query.}
\end{tabular}}
\end{table}

The two largest observed degradations occur in \emph{No Specialist Policy} ($\Delta\text{Pass}_{.50} = -10.0$\,pp, 95\% CI $[-19.0, -1.0]$\,pp; $\Delta\text{Fact} = -6.85$\,pp, CI $[-13.0, -0.7]$\,pp) and \emph{Uniform Text-RAG Baseline} ($\Delta\text{Pass}_{.50} = -9.0$\,pp, CI $[-15.0, -4.0]$\,pp; $\Delta\text{Src.R} = -12.3$\,pp, CI $[-16.5, -8.3]$\,pp). Removing tool isolation causes no change in point estimate ($\Delta\text{Pass}_{.50} = 0.0$\,pp, CI $[-6.0, +6.0]$\,pp) while expanding input tokens by $+1{,}270$. Disabling route repair yields $\Delta\text{Pass}_{.50} = -2.0$\,pp (CI $[-7.0, +3.0]$\,pp), disproportionately harming cross-domain queries ($-15.0$\,pp, dropping from 80.0\% to 65.0\%). Disabling deterministic calculation produces $\Delta\text{Pass}_{.50} = -4.0$\,pp (CI $[-9.0, 0.0]$\,pp) across the benchmark; on the arithmetic tuition subset ($N{=}9$), exact-match answers drop from 8/9 to 3/9 ($-55.6$\,pp exact match pass) due to unconstrained arithmetic hallucination. Collectively, the ablations distinguish two functional roles. First, specialist policies, heterogeneous evidence, and deterministic calculators drive accuracy at the 11-tool scale. Second, tool isolation (S2-A2) acts primarily as an efficiency mechanism ($+1{,}270$ input tokens, $\Delta\text{Pass}_{.50} = 0.0$\,pp); its contribution to accuracy is further examined at larger registry scales in Scenario~3 (Section~\ref{sec:s3_results}). Query family breakdown in Panel~B reveals distinct performance profiles across domain modalities: while cross-domain queries reach $80.0\%$ pass rate, graph-based curriculum queries record $63.2\%$. This gap reflects the structural exactness demanded by curriculum queries, which require multi-property credit aggregations across degree blocks (general, core, elective, and prerequisite chains), where omitting a single credit constraint leads to automated lexical failure. In contrast, cross-domain queries benefit directly from the route-repair mechanism, which resolves domain dispatch conflicts and maintains evidence coverage across disparate policy documents (evidenced by the $15.0$\,pp drop from $80.0\%$ to $65.0\%$ when route repair is disabled).

%------------------------------------------------
\subsection{Scenario 3 --- Tool Ambiguity Stress Test}
\label{sec:s3_results}

Table~\ref{tab:tool_ambiguity} reports the semantic-neighbor experiment across 11--51 tools ($N{=}50$, 3 counterbalanced tool-order blocks, 1,800 primary decisions).

\begin{table}[!htbp]
\centering
\caption{Scenario~3: Semantic-neighbor tool ambiguity ($N=50$, 3 counterbalanced tool-order blocks).}
\label{tab:tool_ambiguity}
{\small\setlength{\tabcolsep}{2.5pt}
\begin{tabular}{llcccccc}
\toprule
\textbf{Reg.} & \textbf{Config.} & \textbf{Vis.} & \textbf{Gold-R} & \textbf{Sel.} & \textbf{E2E} & \textbf{Tok} & \textbf{Lat} \\
\midrule
\multirow{4}{*}{11} & Full-reg single & 11 & -- & \textbf{94.0} & \textbf{86.0} & 1663 & \textbf{1098} \\
 & Top-10 single   & 10 & \textbf{100.0} & 93.3 & 83.3 & 1541 & 1148 \\
 & Routed gen.     & 5.1 & 98.0 & 90.7 & 80.0 & \textbf{1204} & 1889 \\
 & \textbf{Routed spec.} & 5.1 & 98.0 & 93.3 & 85.3 & 1441 & 1983 \\
\midrule
\multirow{4}{*}{31} & Full-reg single & 31 & -- & 91.3 & 84.7 & 2706 & 1157 \\
 & Top-10 single   & 10 & \textbf{98.0} & 90.0 & 81.3 & \textbf{1149} & \textbf{1033} \\
 & Routed gen.     & 9.8 & \textbf{98.0} & 88.0 & 79.3 & 1396 & 1958 \\
 & \textbf{Routed spec.} & 9.8 & \textbf{98.0} & \textbf{93.3} & \textbf{88.7} & 1830 & 2015 \\
\midrule
\multirow{4}{*}{51} & Full-reg single & 51 & -- & 92.7 & 83.3 & 3722 & 1289 \\
 & Top-10 single   & 10 & 96.0 & 86.7 & 78.7 & \textbf{1036} & \textbf{997} \\
 & Routed gen.     & 9.8 & \textbf{98.0} & 91.3 & 82.7 & 1340 & 1907 \\
 & \textbf{Routed spec.} & 9.8 & \textbf{98.0} & \textbf{92.7} & \textbf{88.0} & 1770 & 2021 \\
\bottomrule
\end{tabular}}
\end{table}

As the registry scales from 11 to 51 tools, the Full-Registry Single Agent's E2E success rate drops from 86.0\% to 83.3\%, while its input tokens increase by 124\%. The Global Top-10 changes from 83.3\% to 78.7\% E2E, alongside a Gold Recall change from 100.0\% to 96.0\%. The Routed Specialist records 88.0\% E2E at 51 tools with 9.8 visible tools. At the 51-tool endpoint, Routed Specialist has a +9.3\,pp E2E difference over Top-10 Single Agent (95\% CI $[+2.0, +17.3]$\,pp) and a +4.7\,pp point-estimate difference over Full-Registry Single Agent (95\% CI $[-1.3, +11.3]$\,pp) with 52.4\% fewer input tokens (1{,}770 vs.\ 3{,}722). The relative difference in performance degradation (Specialist vs.\ Top-10) is $+7.3$\,pp (95\% CI $[0.0, +15.3]$\,pp).

Stage-level timing across 1,800 decisions shows that Routed Specialist maintains consistent latency ($2{,}091$--$2{,}116$\,ms) as registries expand from 11 to 51 tools, with supervisor routing contributing a fixed ${\sim}950$\,ms overhead. In contrast, Full-Registry Single Agent gating latency grows from $1{,}098$\,ms to $1{,}272$\,ms as prompt schema tokens scale from $1{,}663$ to $3{,}722$, whereas Routed Specialist bounds schema tokens between $1{,}441$ and $1{,}770$.

%------------------------------------------------
\subsection{Supporting Retrieval Validation}
\label{sec:retrieval_results}

Table~\ref{tab:retrieval_evaluation} validates the retrieval subsystem independently.

\begin{table}[!htbp]
\centering
\caption{Supporting evaluation: Progressive retrieval component stacking ($N=100$).}
\label{tab:retrieval_evaluation}
{\small\setlength{\tabcolsep}{5pt}
\begin{tabular}{lccccc}
\toprule
\textbf{Config.} & \textbf{Hit@1} & \textbf{Hit@3} & \textbf{P@5} & \textbf{R@5} & \textbf{MRR@10} \\
\midrule
BM25 Lexical          & 0.530 & 0.770 & 0.206 & 0.617 & 0.669 \\
Dense Semantic        & 0.460 & 0.650 & 0.176 & 0.524 & 0.560 \\
Hybrid RRF            & 0.520 & 0.790 & 0.224 & 0.662 & 0.673 \\
Hybrid + Reranker     & 0.700 & 0.910 & 0.254 & 0.751 & 0.805 \\
\textbf{Full Retrieval} & \textbf{0.740} & \textbf{0.920} & \textbf{0.302} & \textbf{0.861} & \textbf{0.833} \\
\bottomrule
\end{tabular}}
\end{table}

BM25 outperforms dense search on Hit@1 (0.530 vs.\ 0.460), while Hybrid RRF balances both (0.790 Hit@3). The cross-encoder reranker yields the largest individual gain (+18.0\,pp Hit@1). The full configuration reaches 0.740 Hit@1, 0.861 Recall@5, and 0.833 MRR@10; because it incorporates domain-specific components, the gain cannot be attributed to one component alone.

%------------------------------------------------
\subsection{Failure Analysis}
\label{sec:failure_analysis}

Table~\ref{tab:failures} categorizes error modes across the 58 failure instances ($<50\%$ Fact Coverage) in Scenario~1 (34 for Single Agent, 24 for CTU-Chat), established through manual qualitative review of individual execution traces and intermediate outputs.

\begin{table}[!htbp]
\centering
\caption{Failure taxonomy across architecture configurations ($N=100$). Categories are non-mutually exclusive; a single query may exhibit multiple error modes.}
\label{tab:failures}
{\footnotesize\setlength{\tabcolsep}{2.2pt}
\begin{tabular}{llcc|llcc}
\toprule
\textbf{Category} & \textbf{Failure Mode} & \textbf{S1-A} & \textbf{S1-C} & \textbf{Category} & \textbf{Failure Mode} & \textbf{S1-A} & \textbf{S1-C} \\
\midrule
Routing   & Corrected misroute     & N/A & 1 & Tool       & Invalid arguments     & 4 & 1 \\
          & Unrecoverable misroute & N/A & 3 &            & Wrong-family call     & 3 & 1 \\
Retrieval & Gold source absent     & 3   & 3 &            & Unnecessary call      & 5 & 1 \\
          & Cross-domain omission  & 3   & 2 & Generation & Result ignored        & 2 & 1 \\
          & Reranker demoted       & 7   & 6 &            & Partial fact omission & 22 & 14 \\
\midrule
\multicolumn{4}{l|}{\textbf{Total Failed Queries ($<50\%$ Fact):}} & \multicolumn{4}{l}{\textbf{S1-A (Single): 34} \qquad \textbf{S1-C (CTU): 24}} \\
\bottomrule
\end{tabular}}
\end{table}

CTU-Chat reduces overall failure instances from 34 to 24. Tool errors drop from 12 in Single Agent (4 invalid arguments, 3 wrong-family calls, 5 unnecessary invocations) to 3 in CTU-Chat due to scoped registries. Retrieval omissions (3 gold absent, 7 vs.\ 6 partial demotions) affect both configurations similarly because pre-generation retrieval was matched. In generation, partial fact omissions under complete retrieval ($sr=1.0$) decrease from 22 in Single Agent to 14 in CTU-Chat, reflecting more structured synthesis in specialist prompt contracts.

Trace-level analysis of the three unrecoverable routing failures reveals three modes: (1)~\emph{Repair-induced misroute} (\texttt{HOUT-MHOP-ACAD-04}): rule-based repair overrode a correct \texttt{general} route to \texttt{academic} on curriculum keywords, causing capability refusal. (2)~\emph{Single-owner generation gap} (\texttt{HOUT-XDOM-01}): all gold sources were retrieved ($sr{=}1.0$), but the specialist's scoped prompt refused generating the financial component. (3)~\emph{Single-owner retrieval gap} (\texttt{HOUT-XDOM-16}): the \texttt{scholarship} specialist missed the target document ($sr{=}0.5$). All three illustrate single-owner dispatch limitations for multi-intent queries, differing in where the pipeline breaks: route repair, generation, and retrieval.

